\chapter{Introduction}

\section{Background and Motivation}
FinTech platforms such as digital banking systems, payment applications, lending portals, and investment services depend heavily on recurring customer activity. When customers discontinue a service, revenue drops, customer lifetime value declines, and acquisition cost pressure increases. Churn prediction has therefore become a high-value analytics problem in customer relationship management. Prior literature reports that retaining existing customers is significantly more cost-effective than acquiring new ones, which makes proactive retention a more profitable strategy than reactive replacement \cite{systematicreview,bhattacharjee2026}.

In FinTech settings, churn may appear as account closure, inactivity, reduced transaction behavior, subscription cancellation, or migration to a competing service. Because these behaviors are often preceded by measurable changes in usage and service interaction, machine learning offers a practical way to detect risk early and support targeted business interventions.

\section{Problem Statement}
Many customer analytics systems in FinTech environments rely on static reports, aggregate KPIs, or manually defined rules. Such approaches provide only retrospective visibility and do not reliably identify which individual customers are likely to leave in the near future. The core problem addressed in this project is the absence of a predictive and explainable system that can analyze historical customer data and classify users according to churn risk before the actual churn event occurs.

\section{Objectives}
The objectives of the proposed project are as follows:
\begin{itemize}
    \item To acquire or assume a structured customer dataset appropriate for churn prediction.
    \item To preprocess the dataset by handling missing values, encoding categorical variables, scaling numeric variables, and treating outliers.
    \item To engineer meaningful features related to tenure, engagement, and complaint behavior.
    \item To train multiple machine learning models, namely Logistic Regression, Decision Tree, Random Forest, and XGBoost.
    \item To tune major hyperparameters using cross-validation based search.
    \item To evaluate model performance using accuracy, precision, recall, F1-score, ROC-AUC, and confusion matrix analysis.
    \item To categorize customers into High, Medium, and Low churn risk bands for business use.
    \item To present the full work in a formal academic report with diagrams, tables, and references.
\end{itemize}

\section{Scope and Assumptions}
This report assumes an anonymized customer churn dataset of approximately 10{,}000 records, similar in structure to public churn datasets used in prior studies \cite{malhotra2017}. The dataset source is treated as unspecified because no proprietary institutional dataset was provided. The work focuses on batch model development, offline evaluation, and dashboard-oriented interpretation. Real-time deployment, continuous retraining, and live experimentation are acknowledged as future work rather than core project scope.

The report emphasizes recall and F1-score for the churn class because missing a true churner may be more costly than incorrectly flagging a marginal case. Code listings are intentionally omitted, and the report concentrates on methodology, design, and documented outputs.

\section{Organization of the Report}
The remainder of the report is organized as follows. Chapter 2 presents the literature survey and identifies major research themes. Chapter 3 discusses the existing system, proposed solution, data assumptions, and system requirements. Chapter 4 explains the architecture, DFD, UML models, and deployment design. Chapter 5 describes technologies, preprocessing, feature engineering, model training, and tuning. Chapter 6 covers testing strategy and evaluation metrics. Chapter 7 presents illustrative results and discussion. Chapter 8 concludes the report and outlines future scope.
